\subsection{Exact boundary terms and the reduced scalar}\label{sec:B4T04}
We use Sections~\ref{sec:B4T01}--\ref{sec:B4T03} with $z=\bar p-w$, $u=d-D^Tz$, and
\[
 \mathcal K(a,t)=(a+t)\log\frac{a+t}{a}-t,
 \qquad \mathcal K(a,-a)=a.
\]
We interpret $0\log0$ by continuity.
The parent is stationary for $H_B$, not for $H_B+\ell$.

\paragraph{Action identity without a Taylor expansion in the removed squares.}
Extend $\bar p$ to the parent support by assigning one to $c$ and zero to $F\setminus\{c\}$. The increment preserves all RC constraints, so the linear term from the parent gradient vanishes exactly. The Bregman identity for each summand of $H_B$ gives
\[
 \boxed{\begin{aligned}
 H_{\rm child}(\bar p)-H_{\rm parent}(p)-\log p_c
 ={}&-3\sum_{F\setminus c}(\log q_e+p_e)
 -\sum_{E'}\mathcal K(p_e,z_e)\\
 &+3\sum_{E'}\mathcal K(q_e,-z_e)
 -\sum_{D\cap I}\mathcal K(h_d,u_d).
 \end{aligned}}
\]
The boundary terms can be checked directly. The square $c$ contributes $\log p_c+4q_c$, the other removed squares contribute $-3\log q_e-4p_e$, and the two removed diagonals contribute $-h_{D_+}-h_{D_-}$. We use the exact identity $\sum_{F\setminus c}p_e=4q_c-h_{D_+}-h_{D_-}$.
The four lines on the ordinary board intersect only at $c$.

\paragraph{Local first-order cancellation.}
With $b(p)=(1-4p)/(1-p)^4$, the local factor and the removed block give
\[
 -3\sum_{F\setminus c}(\log q_e+p_e)
 +\tfrac12\sum_F\log b(p_e)+\tfrac12\log\det S_{JJ}=O(k^{-2}).
\]
Indeed, the three leading terms are $\tfrac32\sum_Fp_e^2$, $-3\sum_Fp_e^2$, and $\tfrac32\sum_Fp_e^2$, respectively.
The multiplicity four of $c$ produces a change of only $O(p_c^2)=O(k^{-2})$. The remainders are bounded by $\sum_Fp_e^3=O(k^{-2})$ and $\|S_{JJ}-I\|^2=O(k^{-2})$, since $J$ has size four.

\paragraph{Expansion at the rebalanced point.}
Let $L=H+\ell+\tfrac12\log(2\pi)$ on both connected supports and set
\[
 R_{c0}=L_{\rm child}(\bar p)-L_{\rm parent}(p)-\log p_c.
\]
A Taylor expansion confined to the surviving coordinates and the positive vacancies, together with Section~\ref{sec:B4T01}, gives
\[
 \boxed{R_{c0}=-\tfrac12\left(\sum_{E'}\frac{z_e^2}{p_e}
                  +\sum_{D\cap I}\frac{u_d^2}{h_d}\right)
 +6\sum_{E'}p_ez_e-\tfrac12\operatorname{tr}W_c
 +\tfrac14\operatorname{tr}W_c^2+O(k^{-2}).}
\]
Here we have used $\sum z^2=O(k^{-2})$, $\sum z^2/p=O(k^{-1})$, and $\sum u^2/h=O(k^{-1})$. The cubic remainders are $O(k^{-2})$.
The change in the local factor on the surviving squares satisfies $-\tfrac12\sum_{E'}\log[b(p+z)/b(p)]=6\sum pz+O(k^{-2})$.

We also have $\operatorname{tr}W_c=O(k^{-1})$. Indeed, the diagonal of $\mathcal E$ is $O(k^{-2})$ after correcting the RC masses, $\|\mathcal E\|_\infty=O(k^{-1})$, and $P_c^+=\Pi_I^\perp G_{II}\Pi_I^\perp=I+O_{\max}(k^{-1})$.
Here $\Pi_I$ projects onto the span of $v_I$. The equality follows by solving $S(a,y)=(0,f)$ for $f\perp v_I$ and projecting $y$.
There are $O(k)$ lines, so $\sum_{ij}|\mathcal E_{ij}|=O(1)$.

\paragraph{Reduced scalar.}
Let $e$ be the residual from Section~\ref{sec:B4T02}, let $\bar R$ denote its response, and set $\bar g=\nabla\ell(\bar p)$ and $\delta_c=(R_{\rm parent}\nabla\ell(p))_c/p_c$.
Sections~\ref{sec:B4T02}--\ref{sec:B4T03} and the expansion $\log(1+\delta_c)=\delta_c+O(k^{-2})$ give
\[
 r_c=L_{\rm child}(p^*)-L_{\rm parent}(p)-\log[p_c(1+\delta_c)]
       =\mathcal C_c+O(k^{-2}),
\]
\[
 \boxed{\begin{aligned}
 \mathcal C_c={}&-\tfrac12\left(\sum_{E'}z_e^2/p_e+
                                \sum_{D\cap I}u_d^2/h_d\right)
 +6\sum_{E'}p_ez_e\\
 &-\tfrac12\operatorname{tr}W_c+\tfrac14\operatorname{tr}W_c^2
 +\tfrac12e^T\bar R e+\bar g^T\bar R e-\delta_c.
 \end{aligned}}
\]
All terms are defined using the parent, the rebalancing of Section~\ref{sec:B3T07}, and one Newton step, without solving the child problem anew. These estimates give $\mathcal C_c=O(k^{-1})$. Sections~\ref{sec:B4T06}--\ref{sec:B4T09} prove $\mathcal C_c=O(k^{-2})$ using the resolvent of the unit-diagonal matrix. This overall cancellation is distinct from the local cancellation proved above.
